% TÍTULO — documento de lectura (skill /dossier). Para quién y qué responde.
% Regenerar:  python3 graficos.py && latexmk -lualatex doc.tex
% Todo el contenido de abajo es de ejemplo: se reemplaza entero.
\documentclass[10pt,a4paper]{article}
\usepackage[spanish,es-noshorthands]{babel}
\usepackage{dossier}
\documento{Proyecto · Título corto}{Título completo del documento}{Autor o equipo}

% Identidad del proyecto (opcional): redefinir la paleta con sus tokens, por ejemplo
% \definecolor{acento}{HTML}{22456F}

\begin{document}

% ============================ PÁGINA 1: todo el documento en una página ============================
\portada{Materia o proyecto · equipo · contexto}
  {Título que dice de qué trata}
  {Subtítulo: la pregunta que el documento responde, en una línea}
  {Para quién es y fecha. Cada cifra lleva en gris su fuente.}

\hitos[0.46]{0/Inicio/02-03/hecho, 0.25/Hito 1/06-04/hecho, 0.5/Hito 2/18-05/pendiente, 0.75/Hito 3/15-06/pendiente, 1/Cierre/13-07/pendiente}

\begin{cifras}[4]
  \cifra{Fuentes revisadas}{9}{5 controles cada una}
  \cifra{Decisiones}{12}{9 cerradas · 3 abiertas}
  \cifra{Fuera de plazo}{28,1\,\%}{de 3.265 tickets \fuente{R14}}
  \cifra{Filas del registro}{21.748}{31 columnas \fuente{R02}}
\end{cifras}

\resumen
\textbf{Qué pasó.} La respuesta primero, en dos o tres oraciones con sujeto y verbo. Una cifra
solo si cambia lo que el lector entiende, con su fuente \fuente{R14}.

\textbf{Qué cambió.} Lo que el lector sabía antes y ya no vale, y por qué.

\textbf{Qué sigue.} La próxima decisión o entrega, con fecha. El detalle está en
\ver{sec:uno}.

\ojo[Antes del 18/05]{Lo único que el lector no puede pasar por alto. Como mucho un aviso así por página.}

\indice

% Si el resumen no llena la portada, una pieza de ancho completo la cierra.
\begin{bloque}[Azul oscuro: hecho. Celeste: en curso. Gris: lo que falta.]
\proceso[2]{Relevamiento/A1/hecho, Datos y calidad/A2/hecho, Modelo/A3/actual, Plan y piloto/A4/futuro}
\end{bloque}

% ============================ PÁGINA 2: una página, un mensaje ============================
\newpage
\section{Título que dice el mensaje de la página}\label{sec:uno}

Un párrafo que pone en contexto lo que muestra la figura. Tres a cinco oraciones; lo que se
puede ver en el gráfico no se repite en el texto.

\figura{fig/f1_ejemplo.pdf}{El pie dice qué se ve y cómo leerlo: quitar duplicados mueve poco
la cifra; el mes de medición la mueve mucho \fuente{R14 · R01}.}

\begin{itemize}
\item \textbf{Una idea por punto.} Una oración que la explica.
\item \textbf{Otra idea.} Lo que es propuesta se rotula como tal \chip{propuesta}.
\end{itemize}

\nota[Cómo se midió]{Contexto que ayuda a leer la página pero no urge.}

\subsection{Un gráfico chico junto a su explicación}
% columna: las dos se alinean por arriba; el texto va en bandera y separa los párrafos.
\noindent\begin{columna}{0.5\linewidth}
\figura{fig/f2_categorias.pdf}{2025 quedó 40 clientes por debajo de 2024 \fuente{R03}.}
\end{columna}\hfill
\begin{columna}{0.46\linewidth}
Dos columnas sirven cuando el gráfico es chico y el texto lo explica. El texto dice qué
significa la figura; los valores ya están escritos sobre las barras.

Un segundo párrafo, si hace falta, queda separado como en el resto del documento.
\end{columna}

\begin{tabularx}{\linewidth}{@{}P{3.2cm}LP{2cm}@{}}
\toprule
\enc{Qué} & \enc{Detalle} & \enc{Fuente} \\
\midrule
\textbf{Primero} & una fila por decisión o dato, en pocas palabras & \fuente{R05} \\
\textbf{Segundo} & la tabla ordena lo que la prosa no puede comparar & \fuente{R09} \\
\bottomrule
\end{tabularx}
\medskip

Más información en \enlace{https://ctan.org/pkg/tcolorbox}{la documentación de tcolorbox}.

\anexo{Glosario}
\begin{glosario}
\termino{Fuera de plazo}{Definición copiada de la fuente del proyecto, no escrita de memoria.}
\termino{Ticket duplicado}{Dos registros del mismo pedido y el mismo motivo, abiertos el mismo día.}
\end{glosario}

\end{document}
